\section{Conclusions and open problems}\label{sec:conclusions}

Scott asked whether the Scott--Vogelius--Nitsche method on an inscribed polygon converges with error $\hG^{3/2}+\hO^k$ and, if not, why not. The answer has two halves. The first, announced in the joint letter \cite{CavalcantePadhiLetter} and proved in full here, concerns the method as printed: it converges, but with the expected error only when the pressure is constant along the wall, because without the pressure traction the penalty must balance the wall pressure and the discrete flow leaks through the wall at the rate $(h/\mu)(p-\pbar)$; restoring the traction with its wall mean removed, a device due to Frachon, Nilsson and Zahedi, gives a well-posed, exactly divergence-free method that attains $\hG^{3/2}+\hO^k$ for general data, and the term $\hG^{3/2}$ is sharp. The second half, which is the main new content of this paper, concerns the difficulty that motivated weak imposition in the first place. The zero-gradient lock is a property of the wall mesh: under the Guzm\'an--Scott angle condition strong imposition has the same two-sided rate $\hG^{3/2}+h^k$, and the lock comes from wall vertices in only two triangles, each costing $\hG|\dn u(z)|$, or from nearly singular wall stars. With Scott's normalisation the penalty must grow with the ratio of bulk to boundary mesh size, on every shape-regular mesh, but necessity cannot be read off a single wall star. Finally, the leak is explicit, with constant $1$ in the natural norms, and its $H^1$ size is the energy of a Stokes flow in the whole domain; incompressibility is what keeps the $H^1$ rate at $1$ rather than $\frac12$.

For practitioners the conclusions are simple. If exactly divergence-free velocities are to be combined with Nitsche's method on a polygonal approximation of a curved wall, the pressure traction must be included, with its wall mean removed, and the outer Dirichlet data should be corrected to zero net flux. The penalty should be scaled with the local chord length ($\gamma=\mu\hG/h\ge25$ is safe for $k=4$ on our meshes), and a mesh in which every wall vertex lies in at least three well-shaped triangles should be used, whatever the imposition. The attainable rate is $\hG^{3/2}$; higher rates require curved elements or boundary corrections \cite{NeilanOtus2021,DurstNeilan2024,LiuNeilanOtus2023,ChenLiu2026}.

Several questions remain open.
\begin{itemize}
\item \emph{Constants in the penalty analysis.} Necessity of $\mu\gtrsim\rho$ is proved on every shape-regular mesh, but with astronomically small constants; numerically $\mu^\ast\approx(15\text{--}16)\,h/\delta_1$, with $\delta_1$ the height of the first layer of triangles. An explicit, moderate constant for general meshes is open. The limit $\gamma^\ast_\infty\approx21.41$ of the threshold on our meshes is identified with a periodic cell value only modulo a discrete localisation inequality and a spline cut-off lemma that we have not proved.
\item \emph{The pressure of the corrected method.} Its lower bound needs a hypothesis on a boundary pressure moment; the cell problem that decides this hypothesis rests on a decay property verified only numerically and on a transplant lemma that is not proved. The sharpness of the rate $2$ for the drag of $\CNS$ at fixed penalty is open.
\item \emph{Other variants.} We have analysed the symmetric Nitsche form. The leak comes from the missing pressure term, not from the symmetry of the viscous terms, so we expect it in nonsymmetric variants as well, but we have not analysed these, nor penalty-free variants \cite{Burman2012,BoiveauBurman2016}, the slip-condition method of Gjerde and Scott \cite{GjerdeScott2022}, other divergence-free families, or cut meshes.
\item \emph{Strong imposition on general anisotropic meshes.} Our two-sided needle law assumes that every needle sits in its own pair at the wall; meshes whose needles do not come in isolated wall pairs are not covered.
\item \emph{Three dimensions.} The companion paper \cite{PadhiThreeD} treats inscribed polyhedra; the loss of the chordwise odd-symmetry cancellation changes several arguments there.
\end{itemize}
Finally, although Nitsche's method does not improve the rate on good meshes, it is robust on bad ones: the Nitsche methods were insensitive to locked wall vertices in all our computations. Proving this robustness, on meshes that violate (M2) at the wall, is perhaps the most natural next question.
